\chapter{Divisibility and the Euclidean Algorithm}

\begin{center}\itshape
The oldest non-trivial algorithm still in daily use is Euclid's, from around 300 BC.\\
Its worst case is governed by the Fibonacci numbers, as Lam\'e proved in 1844, in the first analysis of an algorithm's complexity.
\end{center}

\section{Division and ideals of \texorpdfstring{$\Z$}{Z}}

\begin{theorem}[Euclidean division]\label{thm:div}
For $a\in\Z$ and $b\in\N^*$ there is a unique pair $(q,r)\in\Z^2$ with $a=bq+r$ and $0\le r<b$.
\end{theorem}
\begin{proof}
Existence: the set $\{a-bq:q\in\Z\}\cap\N$ is non-empty (take $q=-|a|$), so it has a least element $r=a-bq$. If $r\ge b$ then $r-b=a-b(q+1)$ would be a smaller element. Uniqueness: if $bq+r=bq'+r'$ with $0\le r,r'<b$, then $b\mid r-r'$ and $|r-r'|<b$, so $r=r'$ and $q=q'$.
\end{proof}

\begin{proposition}\label{prop:subgroups}
Every subgroup of $(\Z,+)$ is of the form $d\Z$ for a unique $d\in\N$.
\end{proposition}
\begin{proof}
Let $H\ne\{0\}$ be a subgroup and $d$ its least positive element. Then $d\Z\subset H$. For $a\in H$, divide: $a=dq+r$ with $r=a-dq\in H$ and $0\le r<d$, so $r=0$. Uniqueness: $d$ is the least positive element of $d\Z$.
\end{proof}

\begin{definition}
For $a,b\in\Z$ not both zero, $a\Z+b\Z$ is a subgroup of $\Z$, hence equal to $d\Z$ for a unique $d\ge1$: this $d$ is the \emph{greatest common divisor} $\gcd(a,b)$.
\end{definition}

The name is justified: $d$ divides $a$ and $b$ (they lie in $d\Z$), and any common divisor of $a$ and $b$ divides every element of $a\Z+b\Z$, in particular $d$. Two immediate but fundamental consequences:

\begin{theorem}[B\'ezout]\label{thm:bezout}
There are $u,v\in\Z$ with $au+bv=\gcd(a,b)$. In particular $a$ and $b$ are coprime if and only if $au+bv=1$ for some integers $u,v$.
\end{theorem}

\begin{theorem}[Gauss' lemma]\label{thm:gauss}
If $a\mid bc$ and $\gcd(a,b)=1$, then $a\mid c$. In particular, if a prime $p$ divides $bc$, it divides $b$ or $c$.
\end{theorem}
\begin{proof}
Write $au+bv=1$. Then $c=acu+bcv$, and $a$ divides both terms. For the second statement, if $p\nmid b$ then $\gcd(p,b)=1$.
\end{proof}

\begin{theorem}[fundamental theorem of arithmetic]\label{thm:fta}
Every integer $n\ge2$ is a product of primes, uniquely up to the order of the factors.
\end{theorem}
\begin{proof}
Existence was proved by minimal counterexample in Tome~0, Chapter~1. Uniqueness, by induction on $n$: if $p_1\cdots p_r=q_1\cdots q_s$, then $p_1$ divides the product of the $q_j$, so by Gauss' lemma (iterated) it divides some $q_j$, hence equals it. Cancel and apply the induction hypothesis.
\end{proof}

For a prime $p$ and $n\ne0$, the \emph{$p$-adic valuation} $v_p(n)$ is the exponent of $p$ in $n$. Uniqueness of factorisation gives $v_p(ab)=v_p(a)+v_p(b)$, and
\[
\gcd(a,b)=\prod_pp^{\min(v_p(a),v_p(b))},\qquad\operatorname{lcm}(a,b)=\prod_pp^{\max(v_p(a),v_p(b))},\qquad\gcd(a,b)\operatorname{lcm}(a,b)=|ab|.
\]

\section{The Euclidean algorithm}

Since $\gcd(a,b)=\gcd(b,a-bq)$ (the two pairs generate the same subgroup), we may replace $(a,b)$ by $(b,r)$ where $r$ is the remainder of $a$ divided by $b$. Starting with $r_0=a$, $r_1=b$ and dividing repeatedly,
\[
r_{k-1}=q_kr_k+r_{k+1},\qquad0\le r_{k+1}<r_k,
\]
the remainders decrease strictly, so some $r_{n+1}$ is $0$, and $\gcd(a,b)=r_n$, the last non-zero remainder.

\textbf{Extended version.} Keep track of coefficients with $r_k=au_k+bv_k$: start with $(u_0,v_0)=(1,0)$, $(u_1,v_1)=(0,1)$, and set $u_{k+1}=u_{k-1}-q_ku_k$, $v_{k+1}=v_{k-1}-q_kv_k$. At the end, $au_n+bv_n=\gcd(a,b)$.

\begin{example}
For $(a,b)=(240,46)$: $240=5\cdot46+10$, $46=4\cdot10+6$, $10=1\cdot6+4$, $6=1\cdot4+2$, $4=2\cdot2$. So $\gcd=2$, and running the coefficients gives $240\cdot(-9)+46\cdot47=2$.
\end{example}

\begin{theorem}[Lam\'e]\label{thm:lame}
If the algorithm applied to $a>b\ge1$ performs $n$ divisions, then $b\ge F_{n+1}$ and $a\ge F_{n+2}$, where $(F_k)$ is the Fibonacci sequence.
\end{theorem}
\begin{proof}
The divisions produce $r_0=a>r_1=b>\dots>r_{n-1}>r_n\ge1$ and $r_{n+1}=0$. We show $r_k\ge F_{n+2-k}$ by descending induction on $k$. For $k=n$: $r_n\ge1=F_2$. For $k=n-1$: $r_{n-1}>r_n\ge1$ so $r_{n-1}\ge2=F_3$. Then each $q_k\ge1$, so $r_{k-1}=q_kr_k+r_{k+1}\ge r_k+r_{k+1}\ge F_{n+2-k}+F_{n+1-k}=F_{n+3-k}$.
\end{proof}

So the worst inputs are consecutive Fibonacci numbers, and since $F_{k}\approx\varphi^k/\sqrt5$, the number of steps is $O(\log b)$ (Exercise~\ref{ex:lame5}).

\section{Exercises}

\begin{exercise}[$\star$]\label{ex:bez}
Find all $(x,y)\in\Z^2$ with $91x+35y=7$.
\end{exercise}

\begin{exercise}[$\star$]\label{ex:gcdn}
For $n\in\N$, compute $\gcd\left(n^2+1,(n+1)^2+1\right)$.
\end{exercise}

\begin{exercise}[$\star$]\label{ex:zeros}
\emph{Legendre's formula.} Show that $v_p(n!)=\sum_{k\ge1}\left\lfloor\frac n{p^k}\right\rfloor$, and find the number of zeros at the end of $1000!$.
\end{exercise}

\begin{exercise}[$\star\star$]\label{ex:sqrt}
Show that $\sqrt n$ is rational only when $n\in\N$ is a perfect square. More generally, show that if $x\in\Q$ is a root of a monic polynomial with integer coefficients, then $x\in\Z$.
\end{exercise}

\begin{exercise}[$\star\star$]\label{ex:mersgcd}
Show that $\gcd\left(2^a-1,2^b-1\right)=2^{\gcd(a,b)}-1$ for $a,b\ge1$.
\end{exercise}

\begin{exercise}[$\star\star$]\label{ex:lame5}
Deduce from Lam\'e's theorem that the number of divisions is at most $5$ times the number of decimal digits of $b$.
\emph{Use $F_{k+1}\ge\varphi^{k-1}$ and $\varphi^5>10$.}
\end{exercise}

\begin{exercise}[$\star\star$]\label{ex:frob}
\emph{The coin problem.} Let $a,b\ge1$ be coprime. Show that the largest integer that cannot be written $ax+by$ with $x,y\in\N$ is $ab-a-b$.
\end{exercise}

\begin{exercise}[$\star\star\star$]\label{ex:supercat}
\begin{enumerate}[label=(\alph*)]
\item Show that $\lfloor2x\rfloor+\lfloor2y\rfloor\ge\lfloor x\rfloor+\lfloor y\rfloor+\lfloor x+y\rfloor$ for all real $x,y$.
\item Deduce that $\dfrac{(2m)!\,(2n)!}{m!\,n!\,(m+n)!}$ is an integer for all $m,n\in\N$.
\end{enumerate}
\end{exercise}

\begin{exercise}[$\star\star\star$]\label{ex:n2n}
Show that for every integer $n\ge2$, $n$ does not divide $2^n-1$. \emph{Hint: let $p$ be the smallest prime factor of $n$, and use Exercise~\ref{ex:mersgcd} and Fermat's little theorem.}
\end{exercise}

\begin{beauty}
Every result of this chapter comes from one fact: the subgroups of $\Z$ are the $d\Z$. In Tome~2 this becomes ``$\Z$ is a principal ideal domain'', and the same argument works verbatim for polynomials over a field, Gaussian integers, and every Euclidean ring. B\'ezout's coefficients are how RSA computes private keys (Chapter~3's notebook), and the coin problem with three or more coins, the \emph{Frobenius problem}, has no closed formula and is NP-hard in general.
\end{beauty}

\section{Solutions}

\begin{solution}{ex:bez}
Euclid: $91=2\cdot35+21$, $35=21+14$, $21=14+7$, $14=2\cdot7$, so $\gcd(91,35)=7$, and back-substituting,
$7=21-14=2\cdot21-35=2\cdot91-5\cdot35$. So $(2,-5)$ is a solution. If $(x,y)$ is another, $91(x-2)=-35(y+5)$, so $13(x-2)=-5(y+5)$. By Gauss' lemma $5\mid x-2$: $x=2+5k$, and then $y=-5-13k$. Conversely these all work: the solutions are $(2+5k,-5-13k)$, $k\in\Z$.
\end{solution}

\begin{solution}{ex:gcdn}
Let $d$ be the gcd. It divides $(n+1)^2+1-(n^2+1)=2n+1$. Since $2n+1$ is odd, $d=\gcd(n^2+1,2n+1)$ also divides $4(n^2+1)-(2n+1)(2n-1)=5$. So $d\in\{1,5\}$. If $d=5$ then $5\mid2n+1$, so $n\equiv2\pmod5$. Conversely, if $n=5k+2$ then $n^2+1=25k^2+20k+5$ and $(n+1)^2+1=25k^2+30k+10$ are multiples of $5$. So the gcd is $5$ if $n\equiv2\pmod5$ and $1$ otherwise.
\end{solution}

\begin{solution}{ex:zeros}
$v_p(n!)=\sum_{m=1}^nv_p(m)=\sum_{m=1}^n\sum_{k\ge1}\mathbf 1_{p^k\mid m}=\sum_{k\ge1}\#\{m\le n:p^k\mid m\}=\sum_{k\ge1}\left\lfloor\frac n{p^k}\right\rfloor$.
The number of final zeros is $\min(v_2,v_5)=v_5(1000!)=200+40+8+1=249$.
\end{solution}

\begin{solution}{ex:sqrt}
If $\sqrt n=\frac pq$ then $nq^2=p^2$, so $v_\ell(n)=2v_\ell(p)-2v_\ell(q)$ is even for every prime $\ell$, and $n$ is a square. In general, let $x=\frac pq$ in lowest terms with $q\ge1$ be a root of $X^d+c_{d-1}X^{d-1}+\dots+c_0$. Multiplying by $q^d$: $p^d=-q\left(c_{d-1}p^{d-1}+\dots+c_0q^{d-1}\right)$, so $q\mid p^d$. As $\gcd(q,p)=1$, Gauss' lemma (iterated) gives $q\mid1$, so $x\in\Z$. The case $X^2-n$ recovers the first statement.
\end{solution}

\begin{solution}{ex:mersgcd}
For $a\ge b$, $2^a-1=2^{a-b}\left(2^b-1\right)+\left(2^{a-b}-1\right)$, so $\gcd(2^a-1,2^b-1)=\gcd(2^{a-b}-1,2^b-1)$. The map $(a,b)\mapsto\gcd(2^a-1,2^b-1)$ thus follows the subtractive Euclidean algorithm on the exponents, which ends at $(\gcd(a,b),0)$ and gives $\gcd\left(2^{\gcd(a,b)}-1,0\right)=2^{\gcd(a,b)}-1$.
\end{solution}

\begin{solution}{ex:lame5}
By induction, $F_{k+1}\ge\varphi^{k-1}$: true for $k=1,2$ ($F_2=1$, $F_3=2\ge\varphi$), and $F_{k+2}=F_{k+1}+F_k\ge\varphi^{k-1}+\varphi^{k-2}=\varphi^k$ since $\varphi^2=\varphi+1$. If $b$ has $d$ digits and the algorithm takes $n$ steps, then $10^d>b\ge F_{n+1}\ge\varphi^{n-1}$. As $\varphi^5\approx11.09>10$, we get $\varphi^{n-1}<10^d<\varphi^{5d}$, so $n-1<5d$, that is $n\le5d$.
\end{solution}

\begin{solution}{ex:frob}
If $ax+by=ax'+by'$ then $a(x-x')=b(y'-y)$, so $b\mid x-x'$ by Gauss' lemma: all representations of $m$ in $\Z^2$ are $(x-kb,\,y+ka)$, $k\in\Z$. By B\'ezout there is one, so there is exactly one with $0\le x\le b-1$; call it canonical.
If its $y$ is negative, then any representation with $y+ka\ge0$ has $k\ge1$, hence $x-kb<0$. So $m$ is a non-negative combination if and only if its canonical $y$ is $\ge0$.
For $m=ab-a-b=a(b-1)+b\cdot(-1)$ the canonical $y$ is $-1$: not representable. For $m>ab-a-b$ with canonical $(x,y)$, $by=m-ax\ge m-a(b-1)>-b$, so $y>-1$, that is $y\ge0$: representable.
\end{solution}

\begin{solution}{ex:supercat}
(a) Write $x=a+\alpha$, $y=b+\beta$ with $a,b\in\Z$ and $\alpha,\beta\in[0,1[$. The inequality becomes $\lfloor2\alpha\rfloor+\lfloor2\beta\rfloor\ge\lfloor\alpha+\beta\rfloor$. If $\alpha+\beta<1$ the right side is $0$. Otherwise $\max(\alpha,\beta)\ge\frac12$, so the left side is at least $1=\lfloor\alpha+\beta\rfloor$.

(b) By Legendre's formula, the $p$-adic valuation of the fraction is $\sum_k\left(\left\lfloor\frac{2m}{p^k}\right\rfloor+\left\lfloor\frac{2n}{p^k}\right\rfloor-\left\lfloor\frac m{p^k}\right\rfloor-\left\lfloor\frac n{p^k}\right\rfloor-\left\lfloor\frac{m+n}{p^k}\right\rfloor\right)$, which is $\ge0$ by (a) with $x=\frac m{p^k}$, $y=\frac n{p^k}$. A rational with all valuations $\ge0$ is an integer.
\end{solution}

\begin{solution}{ex:n2n}
Suppose $n\mid2^n-1$; then $n$ is odd, so its smallest prime factor $p$ is odd and $p\mid2^n-1$. By Fermat (Tome~0, Exercise~1.6), $p\mid2^{p-1}-1$. By Exercise~\ref{ex:mersgcd}, $p$ divides $\gcd(2^n-1,2^{p-1}-1)=2^{\gcd(n,p-1)}-1$. But every prime factor of $n$ is $\ge p>p-1$, so $\gcd(n,p-1)=1$, and $p\mid2^1-1=1$: contradiction.
\end{solution}
